% Consolidated for Edward Baker, 4 November 2026.
% Drafting assistant: gpt-6-astra, effort ultra.
% Source: RGD Notes 10--13; Reviews 08 and 10; source ledgers 02--03.
% This is an input fragment, not a standalone manuscript.
% Requires amsmath, amssymb, amsthm, and environments proposition, lemma,
% theorem, corollary, remark.  Labels use the prefix rgd:.
% Bibliography keys used: GarciaGarciaTierz2020Accepted,
% KenanFifteenShifts2026, Odlyzko1987AuthorCopy,
% NISTDLMFPositiveThetaPrimitives.
% The enclosing manuscript should identify Edward Baker as author (or say
% drafted for him) and acknowledge substantial LLM assistance.  The internal
% E04 assessments are not independent human specialist reviews.

\section{Rectangular duality and a finite reciprocal reference}
\label{rgd:sec-reciprocal}

We retain the complete theta source and its normalization,
\begin{align}
 C(z)&=\int_0^\infty\Phi(u)\cosh(u\sqrt z)\,du
      =\frac18\xi\!\left(\frac{1+\sqrt z}{2}\right),
                                                               \label{rgd:C}\\
 C(z)&=\sum_{m\ge0}c_mz^m,\qquad C(1)=b=\frac1{16},
       \qquad c_m=0\quad(m<0),\nonumber\\
 \Phi(u)&=\sum_{n\ge1}
 (2\pi^2n^4e^{9u}-3\pi n^2e^{5u})e^{-\pi n^2e^{4u}}.
                                                               \label{rgd:source}
\end{align}
The square root in \eqref{rgd:C} creates no branch ambiguity, since
both expressions are even in it.  In particular, $c_m>0$ for $m\ge0$.
Put
\[
 \begin{gathered}
 D_{r,k}=\det[c_{k+j-i}]_{0\le i,j<r},\qquad
 E(z)=\frac{C(z)}{c_0}=\sum_{m\ge0}e_mz^m,\\
 H(z)=\frac1{E(-z)}=\sum_{m\ge0}h_mz^m,
 \end{gathered}
\]
where both coefficient sequences are extended by zero at negative indices.
The series for $H$ is initially its Taylor series at zero.

\begin{proposition}[Exact rectangular duality]
\label{rgd:duality}
For every $r,k\ge1$,
\begin{equation}
 D_{r,k}=c_0^r\det[h_{r+j-i}]_{0\le i,j<k}.
 \label{rgd:dual-equation}
\end{equation}
For $k=0$ the left side is $c_0^r$ and the dual determinant is the
empty determinant $1$.  No restriction relating $r$ and $k$ is needed.
\end{proposition}
\begin{proof}
Let $T=[e_{j-i}]_{0\le i,j<r+k}$.  It is upper triangular with diagonal
$1$.  The formal identity $E(z)H(-z)=1$ gives
\[
 (T^{-1})_{ij}=(-1)^{j-i}h_{j-i}.
\]
Take the minor with rows $I=\{0,\ldots,r-1\}$ and columns
$J=\{k,\ldots,k+r-1\}$.  Jacobi's complementary-minor identity expresses
it as $(-1)^{\sum I+\sum J}$ times the inverse minor with rows $J^c$ and
columns $I^c$, since $\det T=1$.  The first sign is $(-1)^{rk}$.
Here $J^c=\{0,\ldots,k-1\}$ and $I^c=\{r,\ldots,r+k-1\}$;
factoring the signs $(-1)^{r+j-i}$ from this inverse minor gives a
second factor $(-1)^{rk}$.  They cancel.  Multiplication by $c_0^r$
proves the result.  All entries at negative indices remain zero.
\end{proof}

This is the rectangular form of the standard elementary/complete
Jacobi--Trudi duality; see, for example,
\cite[\S2.1]{GarciaGarciaTierz2020Accepted}.  We give the finite-matrix
proof to specify the signs and the boundary convention.  The reciprocal
pole comparison below has also been reported in prior work
\cite{KenanFifteenShifts2026}; it is presented here as a method audit,
with no claim of a new fixed-index mechanism or positivity range.\footnote{
The source record attributes the web note to kenan.works, without inferring
a personal byline.  Its primary indexed text reports all-rank coverage
for the first fifteen shifts and the same reciprocal-margin mechanism.
Direct retrieval returned HTTP 502; its supporting notebooks and reported
coverage were not audited.  That report is prior-art attribution, not a
premise of the proofs given here.}

Fix an integer $k\ge1$.  The \emph{finite pole premise} used in this
section is that $C(-z)$ has exactly $k$ zeros in $|z|\le S$, all simple
and positive, with
\begin{equation}
 0<\alpha_1<\cdots<\alpha_k<\rho<S,
 \qquad\text{and no zero on }|z|=S.
 \label{rgd:packet}
\end{equation}
This is a premise about a finite disk.  We do not assume its validity
for arbitrary $k$ or radius.  Define
\begin{equation}
 A_\ell=\frac{c_0}{\alpha_\ell C'(-\alpha_\ell)},\qquad
 H(z)=\sum_{\ell=1}^k\frac{A_\ell}{1-z/\alpha_\ell}+G(z).
 \label{rgd:partial-fractions}
\end{equation}
The complete remainder $G$ is analytic in a neighborhood of $|z|\le S$.
Since $C(0)>0$ and its real zeros in this interval are simple,
successive crossings show
\begin{equation}
 \operatorname{sgn}A_\ell=(-1)^{\ell-1}.
 \label{rgd:residue-sign}
\end{equation}
The actual residue of $H$ at $\alpha_\ell$ is $-A_\ell\alpha_\ell$.

\begin{proposition}[Positive finite-pole margin and complete residual]
\label{rgd:pole-margin}
Assume \eqref{rgd:packet}.  Let $r\ge\max\{1,k-1\}$ and put
\[
 B_{ij}=\sum_{\ell=1}^k A_\ell\alpha_\ell^{-(r+j-i)},
 \quad 0\le i,j<k,\qquad
 \Delta(\alpha)=\prod_{i<j}(\alpha_j-\alpha_i).
\]
Then
\begin{equation}
 d^*_{r,k}:=\det B
 =\Delta(\alpha)^2\prod_{\ell=1}^k
       |A_\ell|\alpha_\ell^{-r-k+1}>0.
 \label{rgd:pole-determinant}
\end{equation}
Choose any $M\ge\max_{|z|=\rho}|G(z)|$, and set
\begin{gather}
 \beta_\ell=\frac{\alpha_\ell}{\rho},\qquad
 a_{\min}=\min_\ell|A_\ell|,\qquad
 u_\ell=\frac{\prod_{j\ne\ell}(1+\beta_j)}
                  {\prod_{j\ne\ell}|\beta_\ell-\beta_j|},
 \quad u_*=\max_\ell u_\ell,                              \label{rgd:pole-constants}\\
 \delta_r=\frac{k^2Mu_*^2}{a_{\min}}\beta_k^r.
 \label{rgd:pole-delta}
\end{gather}
If $\delta_r<1$, then
\begin{equation}
 D_{r,k}\ge c_0^r d^*_{r,k}(1-\delta_r)^k>0.
 \label{rgd:pole-floor}
\end{equation}
In particular, $\delta_r\le\tfrac12$ gives the explicit floor
$2^{-k}c_0^r d^*_{r,k}$.
\end{proposition}
\begin{proof}
The index restriction makes every $r+j-i$ nonnegative.  Factor $B$ into
the Vandermonde matrices $[\alpha_\ell^i]_{i,\ell}$ and
$[\alpha_\ell^{-j}]_{\ell,j}$, separated by
$\operatorname{diag}(A_\ell\alpha_\ell^{-r})$.
The inverse-node Vandermonde has determinant
\[
 (-1)^{k(k-1)/2}\Delta(\alpha)
       \prod_{\ell=1}^k\alpha_\ell^{-(k-1)}.
\]
Its sign cancels the product of the signs in
\eqref{rgd:residue-sign}, proving \eqref{rgd:pole-determinant}.

Write $G(z)=\sum_{m\ge0}g_mz^m$ and
$R_{ij}=g_{r+j-i}$.  Cauchy's coefficient estimate gives
$|g_m|\le M\rho^{-m}$.  Scale each matrix entry $(i,j)$ by
$\rho^{r+j-i}$, and denote the resulting matrices by
$\widetilde B$ and $\widetilde R$.  Each determinant acquires the
positive multiplier $\rho^{rk}$, and
$\|\widetilde R\|_\infty\le kM$.  Moreover,
\[
 \widetilde B
 =U\operatorname{diag}(A_\ell\beta_\ell^{-r})V,
 \qquad U_{i\ell}=\beta_\ell^i,\quad V_{\ell j}=\beta_\ell^{-j}.
\]
The rows of $U^{-1}$ contain the coefficients of the Lagrange
polynomials at the positive nodes $\beta_\ell$.  Their absolute row
sums are exactly $u_\ell$, so $\|U^{-1}\|_\infty=u_*$.
If $J$ reverses the columns, then
\[
 V=\operatorname{diag}(\beta_\ell^{-(k-1)})U^{\mathsf T}J,
 \qquad
 \|V^{-1}\|_\infty
 \le\sum_\ell\beta_\ell^{k-1}u_\ell\le ku_*.
\]
Consequently, for the complete preconditioned remainder,
\[
 \|\widetilde B^{-1}\widetilde R\|_\infty
 \le ku_*\frac{\beta_k^r}{a_{\min}}u_*kM=\delta_r.
\]
When $\delta_r<1$, the real matrices
$I+s\widetilde B^{-1}\widetilde R$ are nonsingular for $0\le s\le1$,
so their determinants remain positive.  Each eigenvalue of
$I+\widetilde B^{-1}\widetilde R$ has modulus at least $1-\delta_r$.
It follows that its determinant is at least $(1-\delta_r)^k$.
Use \eqref{rgd:dual-equation} and the positive determinant scaling to
obtain \eqref{rgd:pole-floor}.  No Hermitianity is assumed.
\end{proof}

For $M>0$, an explicit starting index is
\begin{equation}
 r_0=\max\left\{1,k-1,
 \left\lceil\frac{\log(2k^2Mu_*^2/a_{\min})}
                        {\log(\rho/\alpha_k)}\right\rceil\right\}.
 \label{rgd:pole-start}
\end{equation}
For $M=0$, take $r_0=\max\{1,k-1\}$.  The chosen majorant satisfies
$\delta_{r+1}=\beta_k\delta_r$, and
\[
 \frac{c_0^{r+1}d^*_{r+1,k}}{c_0^rd^*_{r,k}}
   =\frac{c_0}{\prod_{\ell=1}^k\alpha_\ell}.
\]
The first equality concerns the majorant, not an asserted equality for
the actual error norm.  The result is a conditional fixed-$k$ estimate;
it does not supply the finite pole premise as $k$ grows.

\subsection{Why a fixed number of poles cannot supply a growing reference}

\begin{proposition}[Fixed-pole obstruction, including negative indices]
\label{rgd:finite-poles}
Let
\[
 H_N(z)=\sum_{\ell=1}^N\frac{A_\ell}{1-z/\alpha_\ell}
       =\frac{P_{N-1}(z)}{Q_N(z)},\qquad
 Q_N(z)=1+q_1z+\cdots+q_Nz^N,
\]
where $\deg P_{N-1}\le N-1$.  Let $h_m^{[N]}$ be its Taylor
coefficients for $m\ge0$, extended by zero for $m<0$.  Then
\begin{equation}
 \det[h_{r+j-i}^{[N]}]_{0\le i,j<k}=0
       \qquad(r\ge N,\ k>N).
 \label{rgd:fixed-pole-zero}
\end{equation}
This includes $k>r+1$.
\end{proposition}
\begin{proof}
The identity $Q_NH_N=P_{N-1}$ gives
\[
 h_m^{[N]}+q_1h_{m-1}^{[N]}+\cdots+q_Nh_{m-N}^{[N]}=0
       \qquad(m\ge N).
\]
In row $i$ of the last column, $m=r+k-1-i\ge r\ge N$.
The preceding $N$ columns exist because $k>N$.  Thus the last column
plus the stated linear combination of those columns is identically
zero, with the zero negative-index convention unchanged.
\end{proof}

When $r\ge k-1$ one may also factor the whole matrix through an
$N$-dimensional exponential space and conclude rank at most $N$.
That rank argument must not be used for $k>r+1$: the exponential
formula does not extend to the artificially zero negative coefficients.
The recurrence proof establishes \eqref{rgd:fixed-pole-zero} there.
The proposition concerns the reference; residual terms can still make
the complete determinant nonzero.

\subsection{A distinct limitation of the triangle--Harnack allowance}
\label{rgd:harnack-section}

There is also an explicit, but deliberately lossy, way to pay $M$ in
Proposition~\ref{rgd:pole-margin}.  Under \eqref{rgd:packet}, put
\[
 F_0(z)=\frac{C(-z)}{c_0\prod_{\ell=1}^k(1-z/\alpha_\ell)},
 \quad
 L=\frac{C(S)}{c_0\prod_{\ell=1}^k(S/\alpha_\ell-1)},
 \quad U_0\ge\log L.
\]
The function $F_0$ is zero-free on the closed disk, $F_0(0)=1$, and
$|F_0|\le L$ on its boundary.  Thus $L\ge1$.  The analytic logarithm
with $\log F_0(0)=0$, followed by Harnack's inequality for
$U_0-\operatorname{Re}\log F_0$, gives
\begin{equation}
 -\frac{2U_0v}{S-v}\le\operatorname{Re}\log F_0(z)
       \le\frac{2U_0v}{S+v}\qquad(|z|=v<S).
 \label{rgd:harnack-log}
\end{equation}
Let
\begin{gather*}
 p_\ell=\prod_{j\ne\ell}|1-\alpha_\ell/\alpha_j|,
 \quad a^-=\min_\ell A_\ell^-,\\
 A_\ell^-=\frac{e^{-2U_0\alpha_\ell/(S+\alpha_\ell)}}{p_\ell},
 \quad
 A_\ell^+=\frac{e^{2U_0\alpha_\ell/(S-\alpha_\ell)}}{p_\ell},
 \\
 \widehat M=\frac{e^{2U_0\rho/(S-\rho)}}
                      {\prod_\ell(\rho/\alpha_\ell-1)}
       +\sum_\ell\frac{A_\ell^+}{\rho/\alpha_\ell-1},
 \qquad
 \widehat\delta_r=\frac{k^2\widehat M u_*^2}{a^-}\beta_k^r.
\end{gather*}
Indeed, the residue formula and \eqref{rgd:harnack-log} give
$A_\ell^-\le|A_\ell|\le A_\ell^+$; the triangle inequality on
$G=H-\sum A_\ell/(1-z/\alpha_\ell)$ gives $\max_{|z|=\rho}|G(z)|\le\widehat M$, so one may choose
$M=\widehat M$.

\begin{proposition}[Lower bound on this particular allowance]
\label{rgd:harnack-stop}
The specified allowance satisfies
\begin{equation}
 \widehat\delta_r\ge
       4^{k-1}\frac{\beta_k^{r+1}}{1-\beta_k}.
 \label{rgd:harnack-lower}
\end{equation}
Consequently it diverges along any sequence of valid packets with
$k\to\infty$, $r\ge k-1$, $k/r$ bounded below by a positive constant,
and $\beta_k\to1$.
\end{proposition}
\begin{proof}
Interpolate the monic polynomial $(x-\tfrac12)^{k-1}$ at the
$\beta_\ell\in(0,1)$.  Comparing leading coefficients and using
$|\beta_\ell-\tfrac12|\le\tfrac12$ gives
\[
 1\le2^{1-k}\sum_\ell
       \frac1{\prod_{j\ne\ell}|\beta_\ell-\beta_j|}
    \le2^{1-k}ku_*.
\]
Also, the last residue term in $\widehat M$ implies
$\widehat M/a^-\ge1/(\rho/\alpha_k-1)=\beta_k/(1-\beta_k)$.
These inequalities prove \eqref{rgd:harnack-lower}.
For the final assertion, divide its logarithm by $k$.
The contribution $(r+1)\log\beta_k/k$ tends to zero, while
$(k-1)\log4/k$ tends to $\log4$ and
$-\log(1-\beta_k)/k\ge0$.
\end{proof}

For the canonical function, the ordinary zero-counting asymptotic
forces the last condition along any unbounded sequence of valid
packets.  To state this precisely, enumerate the moduli $R_j$ of the
zeros of $C(-z)$ with multiplicity.  A nontrivial zero
$\sigma+i\gamma$ of $\xi$, with $\gamma>0$, maps to
$\alpha=-(2\sigma-1+2i\gamma)^2$, so
$|\alpha|=4\gamma^2+(2\sigma-1)^2$.  The standard zero count
\cite[\S2]{Odlyzko1987AuthorCopy} and $0<\sigma<1$ give
\[
 n(R):=\#\{j:R_j\le R\}\sim\frac{\sqrt R\log R}{8\pi}.
\]
If $R_{j+1}/R_j$ failed to tend to one, two fixed relative interior
radii in a recurring gap would have equal counts but asymptotically
different values under this formula.  Thus $R_{j+1}/R_j\to1$.
A valid packet has $R_k=\alpha_k<\rho<S<R_{k+1}$, whence
$\beta_k\to1$.  The existence of arbitrarily large valid packets is
not assumed; it would in particular impose global reality and
simplicity on the zeros.

This is a failure of the defined allowance, not a lower bound for the
actual remainder.  For example,
$E(z)=\prod_{\ell=1}^k(1+z/\alpha_\ell)$ has exact reciprocal partial
fractions and $G=0$, whereas $\widehat M$ above is positive and can be
large.  Neither Proposition~\ref{rgd:finite-poles} nor
Proposition~\ref{rgd:harnack-stop} excludes every finite-pole argument
or signs a growing-index canonical determinant.

\section{A positive Gaussian reference and its growth obstruction}
\label{rgd:sec-gaussian}

We now consider a reference with a nonzero determinant at every
$r,k$.  Unlike the fixed-pole reference, it needs no finite-zero input.
The primitive for the complete source is
\[
 \begin{gathered}
 p(u)=\sum_{n\ge1}e^{u-\pi n^2e^{4u}},\qquad
 p''-p=8\Phi,\quad p'(0)=-\tfrac12,\\
 c_0=b-\frac{d_0}{8},\qquad d_0=\int_0^\infty p(u)\,du.
 \end{gathered}
\]
These are the canonical Green normalization identities; the derivative
at zero follows from the theta functional equation.  The following
bounds are elementary consequences of the complete sums.

\begin{lemma}[Complete-source envelope]
\label{rgd:envelope}
For $u\ge0$,
\begin{equation}
 0<\Phi(u)\le4\pi^2e^{9u-\pi e^{4u}},\qquad
 \frac1{32}<c_0<\frac1{16}.
 \label{rgd:envelope-equation}
\end{equation}
Writing $a_R=(9+\sqrt R)/4$, for every $R>0$ one has
\begin{equation}
 E(R)<\mathcal F(R):=32\pi^{2-a_R}\Gamma(a_R),
 \qquad \log\mathcal F(R)=O(\sqrt R\log R)\quad(R\ge2)
 \label{rgd:entire-growth}
\end{equation}
as an upper growth estimate.
\end{lemma}
\begin{proof}
Write $v=\pi e^{4u}$, and let $\Phi_n$ be the $n$th summand of
\eqref{rgd:source}.  Every summand is positive, and for $n\ge2$,
\[
 \frac{\Phi_n}{\Phi_1}
 =\frac{n^2(2n^2v-3)}{2v-3}e^{-(n^2-1)v}
 \le2n^4e^{-(n^2-1)v}.
\]
The elementary bounds $n^4\le16^{n-1}$ and
$n^2-4\ge5(n-2)$ show
\[
 \sum_{n\ge2}\frac{\Phi_n}{\Phi_1}
 \le\frac{32e^{-3v}}{1-16e^{-5v}}\le64e^{-3v}<1
       \qquad(v\ge\pi).
\]
Since $\Phi_1\le2\pi^2e^{9u-\pi e^{4u}}$, this proves the envelope.
Also $e^{4u}\ge1+4u$ gives
\[
 0<d_0\le\sum_{n\ge1}\frac{e^{-\pi n^2}}{4\pi n^2-1}
 \le\frac{e^{-\pi}}{(4\pi-1)(1-e^{-3\pi})}<\frac14.
\]
The Green normalization gives the asserted bounds on $c_0$.
For real $R>0$, bound the hyperbolic cosine by $e^{u\sqrt R}$ and
substitute $v=\pi e^{4u}$ in the source integral to obtain
\[
 E(R)\le\frac{\pi^{2-a_R}}{c_0}
       \int_\pi^\infty v^{a_R-1}e^{-v}\,dv
       <32\pi^{2-a_R}\Gamma(a_R).
\]
For completeness, if $N=\lceil a_R\rceil$, then
$v^{a_R-1}\le1+v^N$ for all $v>0$, and hence
$\Gamma(a_R)\le1+N!$.  The estimate $\log N!\le N\log N$ proves
the upper growth assertion without an asymptotic gamma formula.
\end{proof}

Define
\begin{equation}
 t=\log(b/c_0)>0,\qquad C_G(z)=c_0e^{tz},\qquad
 \Phi_G(u)=\frac{c_0}{\sqrt{\pi t}}e^{-u^2/(4t)}.
 \label{rgd:gaussian}
\end{equation}
The elementary Gaussian integral gives the cosh transform $C_G$,
so $C_G(0)=c_0$ and $C_G(1)=b$.  The boundary-compatible primitive
\[
 p_G(u)=8\int_u^\infty\sinh(v-u)\Phi_G(v)\,dv
\]
satisfies $p_G''-p_G=8\Phi_G$ and $p_G'(0)=-8b=-1/2$.
Tonelli's theorem gives
$\int_0^\infty p_G=8\int_0^\infty(\cosh v-1)\Phi_G(v)\,dv
=8(b-c_0)=d_0$.  Thus this comparison retains $C(0)$, $C(1)$, the derivative
$p'(0)$, and the primitive integral.  No equality of $p_G(0)$
and $p(0)$ is asserted.
The complete signed difference
\[
 \Delta E(z)=\frac{C(z)-C_G(z)}{c_0},\qquad
 \Delta E(0)=\Delta E(1)=0
\]
has coefficients
\begin{equation}
 \varepsilon_m=e_m-\frac{t^m}{m!}
 =\frac1{c_0(2m)!}\int_0^\infty
       (\Phi-\Phi_G)(u)u^{2m}\,du\quad(m\ge0).
 \label{rgd:complete-difference}
\end{equation}
They are set to zero for $m<0$.  All source modes and tails remain in
this difference; no bound on its two summands separately is used below.

\begin{proposition}[Gaussian reference margin at every size]
\label{rgd:gaussian-margin}
Let $g_m=t^m/m!$ for $m\ge0$ and $g_m=0$ for $m<0$, and define
$r$-square real matrices
\[
 A=[e_{k+j-i}]_{0\le i,j<r},\qquad
 B=[g_{k+j-i}]_{0\le i,j<r},\qquad K=B^{-1}(A-B).
\]
For every $r,k\ge1$,
\begin{equation}
 \begin{gathered}
 m_{r,k}:=\det B
   =t^{rk}\prod_{j=0}^{r-1}\frac{j!}{(k+j)!}
   \ge\left(\frac{t}{r+k}\right)^{rk}>0,\\
 \frac{D_{r,k}}{c_0^rm_{r,k}}=\det(I+K).
 \end{gathered}
 \label{rgd:gaussian-margin-equation}
\end{equation}
\end{proposition}
\begin{proof}
Factor $t^{k+j}/(k+j)!$ from column $j$ and $t^{-i}$ from row $i$.
The remaining entry is the falling factorial $(k+j)_i$, a monic
polynomial of degree $i$ evaluated at $k+j$.  Its evaluation determinant
is the Vandermonde at $k,k+1,\ldots,k+r-1$, namely
$\prod_{j=0}^{r-1}j!$.  If $i>k+j$, the falling factorial contains a
zero and exactly reproduces the negative-index convention.  The total
power of $t$ is $rk$.  Each quotient $(k+j)!/j!$ consists of $k$
factors at most $r+k$, giving the lower bound.  The final equality
is the exact determinant identity, using $D_{r,k}=c_0^r\det A$.
\end{proof}

The explicit determinant is standard algebra (equivalently a
rectangular Schur specialization); no new determinant identity is
asserted.  It removes the fixed-pole determinant-vanishing obstruction of
Proposition~\ref{rgd:finite-poles}.  Its scale nevertheless prevents a
uniform positive relative slack for the actual complete residual.

\begin{theorem}[Growth obstruction to uniform Gaussian slack]
\label{rgd:gaussian-stop}
For every $r,k\ge1$ and $R>0$,
\begin{align}
 \left|\frac{D_{r,k}}{c_0^rm_{r,k}}\right|
 &\le r^{r/2}E(R)^r\left(\frac{r+k}{tR}\right)^{rk}
 \le\overline U_{r,k}(R),                              \label{rgd:ratio-bound}\\
 \overline U_{r,k}(R)
 &:=r^{r/2}\mathcal F(R)^r
                         \left(\frac{r+k}{tR}\right)^{rk}.
 \nonumber
\end{align}
Fix $M>0$ and set $R_k=k^2/(\log k)^2$ for $k\ge3$.  Then, uniformly
for $1\le r\le Mk$ as $k\to\infty$,
\begin{equation}
 \frac1r\log\overline U_{r,k}(R_k)
 \le-k\log k+2k\log\log k+C_Mk,
 \qquad \overline U_{r,k}(R_k)^{1/r}\longrightarrow0.
 \label{rgd:ratio-decay}
\end{equation}
For every induced matrix norm of the exact residual $K$,
\begin{equation}
 \|K\|\ge1-\overline U_{r,k}(R_k)^{1/r}.
 \label{rgd:norm-stop}
\end{equation}
Thus no bound $\|K\|\le\delta_*<1$ can hold uniformly on this domain.
\end{theorem}
\begin{proof}
Positivity of the coefficients gives $e_mR^m\le E(R)$ for $m\ge0$.
Multiply entry $(i,j)$ of $A$ by $R^{k+j-i}$.  This is an invertible
row and column scaling, with determinant multiplier $R^{rk}$.  All
scaled entries have modulus at most $E(R)$; negative-index entries
remain zero.  Hadamard's inequality gives
\[
 |\det A|\le r^{r/2}E(R)^rR^{-rk}.
\]
Divide by the lower bound for $m_{r,k}$ in
\eqref{rgd:gaussian-margin-equation} and use
Lemma~\ref{rgd:envelope} to get \eqref{rgd:ratio-bound}.
At $R_k$, the upper growth bound gives
$\log\mathcal F(R_k)\le Ck$ for all sufficiently large $k$.
Since $r+k\le(M+1)k$,
\[
 \begin{aligned}
 \frac1r\log\overline U_{r,k}(R_k)
 &=\frac12\log r+\log\mathcal F(R_k)\\
 &\quad+k\bigl(\log(r+k)-\log t-\log R_k\bigr)\\
 &\le-k\log k+2k\log\log k+C_Mk.
 \end{aligned}
\]
All constants are independent of $r$ in the stated domain.  This
proves \eqref{rgd:ratio-decay}, and also uniform convergence of
$\overline U_{r,k}(R_k)$ itself to zero because $r\ge1$.
If $\|K\|<1$, every eigenvalue of $I+K$ has modulus at least
$1-\|K\|$.  The determinant identity consequently gives
\[
 (1-\|K\|)^r\le|\det(I+K)|
      \le\overline U_{r,k}(R_k),
\]
which proves \eqref{rgd:norm-stop}.  If $\|K\|\ge1$, that inequality
is immediate.
\end{proof}

The theorem includes fixed $r$, $r=k$, and any fixed proportional
ratio with $k/r\ge1/M$.  It does not cover fixed $k$ with $r\to\infty$.
It is a lower bound $\|K\|\ge1-o(1)$, not a claim that $\|K\|$ tends
to one.  If $A,B$ undergo common invertible scalings $A\mapsto LAR$
and $B\mapsto LBR$, then $K\mapsto R^{-1}KR$.  Norm values may change,
but the determinant-based obstruction remains valid for every induced
norm of the transformed residual.

There is a corresponding exact scalar obstruction:
\begin{equation}
 \frac{\det A-m_{r,k}}{m_{r,k}}\longrightarrow-1,
 \qquad
 \frac{|\det A-m_{r,k}|}{m_{r,k}}\longrightarrow1
 \label{rgd:scalar-stop}
\end{equation}
uniformly on the same domain.  It excludes a fixed positive relative
slack, not every tolerance approaching one.  Once $|\det A|<m_{r,k}$,
the bare inequality $|\det A-m_{r,k}|<m_{r,k}$ is equivalent to
$\det A>0$.  A proof using such a shrinking tolerance must therefore
supply an independently established signed estimate; merely imposing
the inequality restates the target.

The scale mismatch can also be made explicit.  If a positive moving
slope $t_{r,k}$ in the exponential family permitted a fixed norm bound
$\delta<1$, the same argument would require
\begin{equation}
 t_{r,k}\le
 \frac{r+k}{R_k(1-\delta)^{1/k}}
       r^{1/(2k)}E(R_k)^{1/k}
 =O_{M,\delta}\!\left(\frac{\log^2 k}{k}\right).
 \label{rgd:moving-slope}
\end{equation}
This is a necessary condition only.  Matching both endpoint values
fixes $t=\log(b/c_0)>0$, so it cannot meet this vanishing scale in the
same single-exponential family.

\subsection{The endpoint-preserving quadratic adjustment}

\begin{proposition}[Failure of universal positivity for a nonzero correction]
\label{rgd:quadratic}
For $q\in\mathbb R$, put
\[
 E_2(z)=\exp\bigl((t+q/2)z-qz^2/2\bigr),\qquad C_2=c_0E_2.
\]
Then $C_2(0)=c_0$ and $C_2(1)=b$.  For every $q\ne0$, some
rectangular coefficient determinant of $C_2$ is negative.
\end{proposition}
\begin{proof}
If $q>0$, then $E_2(x)\to0$ as $x\to+\infty$.  An entire series
with constant coefficient one and all coefficients nonnegative is
at least one on that axis.  Thus some coefficient $e_j^{[2]}$ is
negative, giving $D_{1,j}^{[2]}=c_0e_j^{[2]}<0$.
If $q<0$, the dual series
\[
 H_2(z)=\frac1{E_2(-z)}
       =\exp\bigl((t+q/2)z+qz^2/2\bigr)
\]
tends to zero on the positive axis, so some $h_j^{[2]}<0$.
Proposition~\ref{rgd:duality}, which is purely algebraic, gives
$D_{j,1}^{[2]}=c_0^jh_j^{[2]}<0$.
\end{proof}

Matching the first coefficient as well would select $q=2(e_1-t)$;
no sign or nonzero assertion about that canonical number is needed.
For $q=0$, Theorem~\ref{rgd:gaussian-stop} applies.  The nonzero
quadratic result concerns positivity at \emph{all} indices and uses
edge-family witnesses; it does not exclude every restricted
proportional-index subdomain or other reference families.

\subsection{A positive control and the extent of the conclusions}

The Gaussian obstruction is compatible with strict positivity of every
target rectangle.  Indeed, for $\ell>0$ take
\[
 C_*(z)=\frac{\cosh(\ell\sqrt z)}{16\cosh\ell},\qquad
 E_*(z)=\cosh(\ell\sqrt z).
\]
The classical product \cite[Eq.~4.36.2]{NISTDLMFPositiveThetaPrimitives}
gives $E_*(z)=\prod_{j\ge1}(1+x_jz)$, where
$x_j=4\ell^2/((2j-1)^2\pi^2)>0$ and $\sum_jx_j<\infty$.
For a finite product, its upper Toeplitz matrix in dimension $r+k$ is
the product of upper bidiagonal matrices $I+x_jS$, with $S$ the upper
shift.  Each factor has nonnegative minors.  In the Cauchy--Binet
expansion for rows $0,\ldots,r-1$ and columns $k,\ldots,k+r-1$,
move the full consecutive block by one at each of the first $k$
factors and leave it fixed thereafter.  This supplies the positive
term $(x_1\cdots x_k)^r$.  Local convergence of the product preserves
that lower bound in every fixed rectangle.  Thus all these target
determinants are strictly positive, while the same proof of
Theorem~\ref{rgd:gaussian-stop}, using
$\log E_*(R)\le\ell\sqrt R$, gives the Gaussian fixed-slack obstruction.
This is an atomic source control, not the canonical smooth theta source
or a certified positive-time collision.

The conclusions here are deliberately separate.  A fixed finite pole
reference eventually has a zero determinant; the specified
triangle--Harnack allowance can diverge despite a zero actual error;
and the endpoint-matched Gaussian reference has a positive margin at
every size but cannot support uniform positive relative slack on the
stated growing-index domain.  None of these facts makes a canonical
determinant negative, supplies a new theta positivity range, or proves
an all-index induction.  Continuation would require an appropriately
scaled reference with a paid complete residual, or an independently
controlled signed comparison.  No such estimate is asserted here.

\paragraph{Assessment provenance.}
The reciprocal results were accepted within their recorded scope in
internal assessment RGD-E04 r5, and the Gaussian and quadratic results
were accepted within their recorded scope in RGD-E04 r6, without a
mathematical correction.  Those assessments used a different LLM model
with shared source, algebra, and envelope premises; they were not
independent human specialist reviews.  The present consolidation adds
no numerical evidence, certificate replay, new positivity coverage, or
priority claim.  The direct algebraic proofs above are the mathematical
support for the statements, not the assessment labels.
